\clearpage
\VerdierPhysicalPageAnchor{0172}%
\begingroup\small\setlength{\parskip}{0pt}\setlength{\emergencystretch}{1.5em}
\VerdierSemanticAnchor{ch4-4-4}
\medskip\noindent\textbf{4.4. Convergence.}

\VerdierSemanticAnchor{ch4-4-4-1}
\medskip\noindent\textbf{4.4.1.}\quad
For a detailed study of the convergence of spectral sequences, we refer to
\emph{loc. cit.} We shall use only the following criterion:

\VerdierSemanticAnchor{ch4-4-4-2}
\medskip\noindent\textbf{Definition 4.4.2.}\quad
A spectral object $(H^n(p,q),\delta)$ with values in an abelian category
$\mathcal A$ is said to be \emph{stationary} if, for every integer $n$,
there exist two integers $u(n),v(n)\in\mathbb Z$ such that:
\[
\begin{aligned}
H^n(p,q)&=0,&&u(n)\leq p\leq q\leq+\infty,\\
H^n(p,q)&=0,&&-\infty\leq p\leq q\leq v(n).
\end{aligned}
\]

\VerdierSemanticAnchor{ch4-4-4-3}
\medskip\noindent\textbf{Theorem 4.4.3.}\quad
\textit{Let $(H^n(p,q),\delta)$ be a stationary spectral object.}

\textit{a) For every $n$, the filtration
$F^pH^n(-\infty,+\infty)$, $p\in\widetilde{\mathbb Z}$, is stationary;
i.e. there exist two integers $n_1\leq n_2$ in $\mathbb Z$ such that:}
\[
F^{n_1}H^n(-\infty,+\infty)=0
\quad\text{and}\quad
F^{n_2}H^n(-\infty,+\infty)=H^n(-\infty,+\infty).
\]

\textit{b) The differentials leaving (respectively, entering)
$E\!\left[\begin{smallmatrix}n\\r,p,q,s\end{smallmatrix}\right]$ are zero
if $r\leq v(n+1)$ (respectively, $u(n-1)\leq s$), and moreover:}

\textit{c) For every $n$ and every pair $p\leq q$, the objects
$E\!\left[\begin{smallmatrix}n\\r,p,q,s\end{smallmatrix}\right]$ such that}
\[
-\infty\leq r\leq v(n+1)
\quad\text{and}\quad
u(n-1)\leq s\leq+\infty
\]
\textit{are canonically isomorphic.}

This theorem follows immediately from the definitions.

\VerdierSemanticAnchor{ch4-4-4-3-convergence}
A stationary spectral object gives a strongly convergent spectral sequence
in the sense of \emph{loc. cit.} We shall say that a spectral sequence
associated with a spectral object is \emph{stationary}, or, by abuse of
language, \emph{convergent}, when that spectral object is stationary. Using
one or the other of the notations in (4.3.3), we shall write:
\[
\mathrm I_r^{p,q}\Longrightarrow H^{p+q}(-\infty,+\infty),
\qquad
\mathrm{II}_r^{p,q}\Longrightarrow H^{p+q}(-\infty,+\infty)
\]
to indicate that a spectral sequence converges. The collection of terms
$H^{p+q}(-\infty,+\infty)$ will then be called the \emph{abutment} of the
spectral sequence.
\endgroup
